% 两变量成对二值函数：phi_1=(0,1)，phi_2=(2,0)，phi_{1,2}(00,01,10,11)=(0,3,5,0)。
% J = phi_1(x1)+phi_2(x2)+phi_{1,2}(x1,x2)：J(00)=2，J(01)=3，J(10)=8，J(11)=1。
% 逐点选择各自一元项的最小者得 (0,1)，J=3；整体最小为 (1,1)，J=1。
\begin{tikzpicture}[x=1mm,y=1mm,line width=.6pt,font=\figurefont\fontsize{8}{10}\selectfont,
  cell/.style={draw=FigureStroke,minimum width=30mm,minimum height=12mm,align=center,line width=.8pt,inner sep=1pt}]
  \node[cell,fill=white] at (15,18) {$0+2+0=2$};
  \node[cell,fill=FigureYellowFill] at (45,18) {$0+0+3=3$\\[-1pt]{\FigureNoteFont\fontsize{7.5}{9}\selectfont 逐点选择}};
  \node[cell,fill=white] at (15,6) {$1+2+5=8$};
  \node[cell,fill=FigureBlueFill] at (45,6) {$1+0+0=1$\\[-1pt]{\FigureNoteFont\fontsize{7.5}{9}\selectfont 整体最小}};
  \node[FigureInk] at (15,27) {$x_2=0$};
  \node[FigureInk] at (45,27) {$x_2=1$};
  \node[FigureInk,anchor=east] at (-1,18) {$x_1=0$};
  \node[FigureInk,anchor=east] at (-1,6) {$x_1=1$};
  \node[FigureMuted,anchor=north,font=\FigureNoteFont\fontsize{7.5}{9}\selectfont] at (30,-1.5) {每格：$\phi_1(x_1)+\phi_2(x_2)+\phi_{1,2}(x_1,x_2)=J(x_1,x_2)$};
\end{tikzpicture}
